\scanpage{25}
En effet, soit $\psi$ une autre polarisation, donc $\widetilde\psi=\widetilde\varphi\circ u$, où $u\in\operatorname{End}(A)\otimes_{\mathbb Z}\mathbb Q$, \deleted{le} et la condition $\widetilde\psi'=-\widetilde\psi$ (jointe à $\widetilde\varphi'=-\widetilde\varphi$) donne
\[
\begin{array}{ccc}
u'\widetilde\varphi'=-\widetilde\varphi u & \text{i.e.} & u'=\widetilde\varphi u\widetilde\varphi^{-1}.\\
\Vert &&\\
-u'\widetilde\varphi &&
\end{array}
\]
On veut prouver [posant $\delta(\widetilde\varphi)=\sqrt{\deg\widetilde\varphi}$, $\alpha_p(\widetilde\varphi)=\Delta(\widetilde\varphi)$]
\[
\begin{array}{ccc}
\Delta(\widetilde\psi)/\delta(\unsure{\widetilde\varphi|\widetilde\psi})&=&\Delta(\widetilde\varphi)/\delta(\widetilde\varphi)\\
\Vert&&\\
\Delta(\widetilde\varphi)\cdot\Delta(u)/\delta(\widetilde\varphi)\delta(u)&&
\end{array}
\]
qui équivaut à: $\Delta(u)/\delta(u)=1\deleted{\unsure{2}}$. Or par \unsure{lemme} on a
\[
\begin{array}{c}
\deleted{\Delta}\;\Delta(u)\Delta(u^*)=\delta(u)^2\\
\hfill\Vert\\
\hfill\deg u
\end{array}
\]
d'autre part, comme $u^*=\widetilde\varphi u\widetilde\varphi^{\deleted{\unsure{*}}\inserted{-1}}$, on a ici $\Delta(u^*)=\Delta(u)$, donc \deleted{$\unsure{\Delta(u)}$}
\[
\Delta(u)^2=\delta(u)^2
\]
d'où $\Delta(u)=\pm\delta(u)$. Il faut encore avoir le signe, il faut \deleted{\unsure{pour dire}} donc sur $\operatorname{End}(A)\otimes_{\mathbb Z}\mathbb Q$, muni de l'involution $u\unsure{\longmapsto|\sim}u'=\widetilde\varphi u\widetilde\varphi^{-1}$, prouver que $u>0$ (i.e. \deleted{\unsure{réellement}} $u=\sum v_iv_i'$) \deleted{$u=u'$} implique $\Delta(u)={\text{\uuline{\uline{$+$}}}}\delta(u)$ [NB. $\delta(u)$ est bien défini, car $\deg u$ est bien un carré parfait si $u=u'$]
\deleted{Or l'identité $\Delta(u)^2=\delta(u)$ est une identité dans}
\mspage{24}
